% Figure: a piecewise assembly. Three panels: Heaviside step,
% signum function, and a piecewise linear ramp-saturate (a soft
% saturator used in controller design).
\begin{figure}[htbp]
\centering
\resizebox{\linewidth}{!}{%
\begin{tikzpicture}[scale=0.92,
  ax/.style={->, draw=engblack, thick},
  tick/.style={draw=engblack!60, thin},
  curve/.style={draw=engblue, very thick},
  hole/.style={draw=engblue, thick, fill=white},
  filled/.style={fill=engblue, draw=engblue, thick},
  ann/.style={font=\footnotesize, align=center},
]

% --- Left: Heaviside step ---
\begin{scope}[xshift=0cm]
  \draw[ax] (-2.5, 0) -- (2.5, 0) node[right, ann] {$x$};
  \draw[ax] (0, -0.5) -- (0, 2.0) node[above, ann] {$H(x)$};
  \draw[tick] (-2, -0.05) -- (-2, 0.05);
  \draw[tick] (2, -0.05) -- (2, 0.05);
  \draw[tick] (-0.05, 1) -- (0.05, 1);
  \node[ann, left] at (-0.1, 1) {$1$};
  \draw[curve] (-2.4, 0) -- (-0.02, 0);
  \draw[curve] (0.02, 1) -- (2.4, 1);
  \node[hole, circle, inner sep=1.3pt] at (0, 0) {};
  \node[filled, circle, inner sep=1.3pt] at (0, 1) {};
  \node[ann, anchor=north] at (0, -0.7) {Heaviside step $H(x)$};
\end{scope}

% --- Middle: signum ---
\begin{scope}[xshift=6cm]
  \draw[ax] (-2.5, 0) -- (2.5, 0) node[right, ann] {$x$};
  \draw[ax] (0, -1.5) -- (0, 1.7) node[above, ann] {$\sgn(x)$};
  \draw[tick] (-2, -0.05) -- (-2, 0.05);
  \draw[tick] (2, -0.05) -- (2, 0.05);
  \draw[tick] (-0.05, 1) -- (0.05, 1);
  \draw[tick] (-0.05, -1) -- (0.05, -1);
  \node[ann, anchor=west] at (0.1, 1.0) {$+1$};
  \node[ann, anchor=west] at (0.1, -1.0) {$-1$};
  \draw[curve] (-2.4, -1) -- (-0.02, -1);
  \draw[curve] (0.02, 1) -- (2.4, 1);
  \node[filled, circle, inner sep=1.3pt] at (0, 0) {};
  \node[hole, circle, inner sep=1.3pt] at (0, 1) {};
  \node[hole, circle, inner sep=1.3pt] at (0, -1) {};
  \node[ann, anchor=north] at (0, -2.0) {signum $\sgn(x)$};
\end{scope}

% --- Right: ramp-saturate (deadband + soft saturator) ---
\begin{scope}[xshift=12cm]
  \draw[ax] (-2.5, 0) -- (2.5, 0) node[right, ann] {$x$};
  \draw[ax] (0, -1.5) -- (0, 1.7) node[above, ann] {$y$};
  \draw[tick] (1, -0.05) -- (1, 0.05);
  \draw[tick] (-1, -0.05) -- (-1, 0.05);
  \draw[tick] (-0.05, 1) -- (0.05, 1);
  \draw[tick] (-0.05, -1) -- (0.05, -1);
  % Saturate: y = -1 for x < -1; y = x for -1 <= x <= 1; y = 1 for x > 1
  \draw[curve] (-2.4, -1) -- (-1, -1) -- (1, 1) -- (2.4, 1);
  \node[ann, anchor=north] at (0, -2.0)
    {soft saturator\\(piecewise linear)};
\end{scope}

\end{tikzpicture}%
}%
\caption[Piecewise definitions assembled from elementary pieces]{Three
piecewise functions the working engineer meets daily. The Heaviside
step $H(x)$, by convention $0$ on the left and $1$ on the right
(value at zero a matter of convention; signal-processing conventions
choose $1/2$ for Fourier work and $1$ for control work). The signum
function $\sgn(x)$, the odd cousin of the step. The soft saturator,
a piecewise linear function that ramps through a central window
$[-1, 1]$ and clips outside it, the canonical model of an op-amp
output stage hitting the supply rails or a controller's actuator
limit. The filled circles mark the value the function actually takes
at a discontinuity; the open circles mark the limit the function
does not take. Sloppy notation about which value the function takes
at the discontinuity is a recurring source of error in published
control diagrams.}
\label{fig:vol02:ch02:piecewise-step}
\end{figure}
